Considereu la funció polinòmica $f(x)=3x^{13}+5x^3+2$.

\begin{apartats}

\apartat{1,25}
Justifiqueu que la seva gràfica talla l'eix de les abscisses en un punt de l'interval $[-2,0]$. Doneu un
interval de longitud 0,5 on es trobi aquest punt de tall.

\begin{solucio}
La funció $f(x)$ és polinòmica i, per tant, contínua en tots els reals. Com que
\[
f(-2)=3\cdot(-2)^{13}+5\cdot(-2)^3+2<0,\qquad f(0)=3\cdot(0)^{13}+5\cdot(0)^3+2>0,
\]
podem aplicar el teorema de Bolzano, i la funció $f(x)$ s'anul·la forçosament en algun punt de l'interval
$[-2,0]$. Per trobar un interval de longitud 0,5, només cal temptejar uns quants intervals més petits. Com que
$f(-1)=3\cdot(-1)^{13}+5\cdot(-1)^3+2=-6<0$, l'arrel pertany a l'interval $[-1,0]$. I, finalment, com que
\[
f(-0{,}5)=3\cdot(-0{,}5)^{13}+5\cdot(-0{,}5)^3+2=1{,}374\ldots>0,
\]
l'arrel pertany a l'interval $[-1,-0{,}5]$.

\textit{Pauta oficial:} 0,25 per justificar que es pot aplicar el teorema de Bolzano, 0,5 per veure que hi ha
una arrel a l'interval $[-2,0]$ i 0,5 per afinar fins a un interval de cinc dècimes d'amplada.
\end{solucio}

\apartat{1,25}
Estudieu les zones de creixement i de decreixement, i els màxims i els mínims de $y=f(x)$. Quants punts de
tall té exactament la gràfica d'aquesta funció amb l'eix de les abscisses? Justifiqueu la resposta.

\begin{solucio}
Per analitzar el creixement i decreixement de la funció $f(x)$, calculem la seva derivada,
$f'(x)=39x^{12}+15x^2$: com que tots els exponents són parells, $f'(x)$ mai és negativa i, per tant, la funció
$f$ mai decreix. En particular, no té cap màxim ni mínim. D'altra banda, la funció, que és polinòmica, no té
asímptotes ni discontinuïtats. En definitiva, la seva gràfica només pot tenir un únic punt de tall amb l'eix
d'abscisses, que es troba, com hem dit, a l'interval $[-1,-0{,}5]$.

\textit{Pauta oficial:} 0,25 pel càlcul de la derivada, 0,25 per deduir que es tracta d'una funció creixent,
0,25 per dir que no té màxims ni mínims, i 0,5 per justificar que hi ha una única arrel.
\end{solucio}

\end{apartats}
